\section*{Appendix B: Conductor sheaf and Grothendieck duality.}
\addcontentsline{toc}{section}{Appendix B: Conductor sheaf and Grothendieck duality}
\renewcommand{\thesection}{B}

In this appendix we detail the coherent sheaf duality theory on the non-normal central fiber $W_0$.

\subsection*{B.1 Grothendieck duality on Cohen--Macaulay surfaces}

Let $W_0$ be the reduced normal crossings central fiber from Section~4, with normalization $\nu \colon \widetilde W_0 \to W_0$ and singular locus $D = \mathrm{Sing}(W_0)$.
The conductor ideal sheaf is defined by:
\[
\mathscr{C}_{W_0} := \operatorname{\mathcal{H}\mathit{om}}_{\cO_{W_0}}(\nu_* \cO_{\widetilde W_0}, \, \cO_{W_0}) \subset \cO_{W_0}.
\]

\begin{theorem}[Conductor Duality Formula]\label{thm:conductor-duality}
\begin{enumerate}[label=(\alph*)]
    \item The dualizing sheaf $\omega_{W_0}$ is an invertible $\cO_{W_0}$-module satisfying:
    \[
    \nu^* \omega_{W_0} \cong \omega_{\widetilde W_0}(D),
    \]
    where $D \subset \widetilde W_0$ is the preimage of the double curve cycle.
    \item For any coherent sheaf $\cF$ on $W_0$, Grothendieck duality yields:
    \[
    \operatorname{Ext}^i_{W_0}(\cF, \, \omega_{W_0}) \cong H^{2-i}(W_0, \, \cF)^\vee.
    \]
    \item Taking $\cF = TX|_{W_0} \otimes L$, the non-vanishing of the conductor section $s = df \otimes e$ establishes:
    \[
    H^2(W_0, \, TX|_{W_0} \otimes L) \ne 0 \quad \text{for every } L \in \Pic(X),
    \]
    providing a complete, first-principles verification of Theorem~\ref{thm:cdp-failure}.
\end{enumerate}
\end{theorem}
